\chapter{Software - Excel}

\begin{outcome}
\begin{itemize}
\item Learn relevant Excel functions
\item Set up linear programs in Excel
\item Solve linear programs in Excel
\end{itemize}
\end{outcome}

\section{Using Excel Solver}
\label{sec:excel-solver}

Excel Solver is a powerful tool for solving optimization problems with linear and nonlinear constraints. Before diving into using Solver, it is essential to understand some basic Excel functions and techniques that are commonly used when setting up optimization problems. These functions help organize data, compute intermediate values, and define objective functions and constraints.

\subsection{Useful Excel Functions and Techniques}

\paragraph{Using Formulas in Cells}
In Excel, formulas begin with an equals sign (\texttt{=}). For example, typing \texttt{=A1+B1} in a cell will calculate the sum of the values in cells \texttt{A1} and \texttt{B1}.

\paragraph{\texttt{SUM()}}
The \texttt{SUM()} function adds a range of numbers. For example, \texttt{=SUM(A1:A10)} calculates the sum of all values in the range \texttt{A1} to \texttt{A10}.

\paragraph{\texttt{SUMPRODUCT()}}
The \texttt{SUMPRODUCT()} function multiplies corresponding elements in two or more arrays and returns the sum of those products. This is particularly useful for computing dot products or weighted sums. For example, \texttt{=SUMPRODUCT(A1:A10, B1:B10)} computes \( \sum_{i=1}^{10} A_i B_i \).

\paragraph{\texttt{COUNT()}}
The \texttt{COUNT()} function counts the number of numeric entries in a range. For example, \texttt{=COUNT(A1:A10)} returns the number of numeric cells in the range \texttt{A1} to \texttt{A10}.

\paragraph{\texttt{COUNTIF()}}
The \texttt{COUNTIF()} function counts the number of cells in a range that meet a specific condition. For example, \texttt{=COUNTIF(A1:A10, ">5")} counts the number of cells in \texttt{A1:A10} with values greater than 5.

\paragraph{Duplicating Formulas Across Rows or Columns}
To apply a formula to multiple cells, you can drag the fill handle (a small square at the bottom-right corner of the selected cell) across the desired range. Excel automatically adjusts cell references based on the relative position (relative referencing). To keep a reference constant, use the dollar sign (\texttt{\$}) in the cell reference (absolute referencing). For example, \texttt{\$A\$1} always refers to cell \texttt{A1}.

\paragraph{Conditional Formatting}
Conditional formatting allows you to highlight cells that meet specific criteria. For example, you can format cells with values above a threshold in bold or a particular color to identify important data visually.

\paragraph{Data Validation}
Data validation helps restrict the type of data entered in a cell. For example, you can set a rule that only allows numbers within a specific range or text from a predefined list.

\paragraph{Naming Ranges}
Naming ranges makes formulas easier to read and maintain. Instead of \texttt{=SUM(A1:A10)}, you can name the range \texttt{"Sales"} and write \texttt{=SUM(Sales)}.

\paragraph{Using Logical Functions}

Logical functions like \texttt{IF()}, \texttt{AND()}, and \texttt{OR()} are helpful for decision-making within formulas. For example, \texttt{=IF(A1>10, "High", "Low")} returns ``High'' if \texttt{A1} is greater than 10, and ``Low'' otherwise.

With these foundational tools, setting up optimization problems becomes much more manageable. In the next section, we will explore how to use these functions in combination with Excel Solver to define and solve optimization models.

\subsection{How to Use Excel Solver}

Excel Solver is a powerful add-in that allows users to define and solve optimization problems. This section provides a step-by-step guide on how to effectively layout data and use Solver for optimization.

\paragraph{Step 1: Layout the Data}
Organizing your data clearly is crucial for using Solver effectively. Use color coding to distinguish between different types of cells:
\begin{itemize}
    \item \textbf{Data:} Cells containing fixed input values, such as coefficients or constants, should be highlighted in \textcolor{blue!50!white}{light blue}.
    \item \textbf{Variables:} Cells representing decision variables should be highlighted in \textcolor{yellow}{yellow}.
        \item \textbf{Formulas:} Cells with formulas, such as the objective function or summarizing constraints values should be highlighted in \textcolor{green!50!white}{light green}.
\end{itemize}
Other color combinations are also acceptable, as long as they consistently distinguish between these components.

{\centering
\altincludegraphics[width = 0.85\textwidth]{Figures/excel-solver-layout}{Excel spreadsheet showing the recommended Solver layout: data cells (coefficients, right-hand sides) shaded light blue, decision-variable cells shaded yellow, and formula cells (objective and constraint sums) shaded light green.}\par}

\paragraph{Step 2: Define the Objective Function}
Identify the cell that represents the objective function. This cell should contain a formula that calculates the value to be maximized or minimized, such as total profit or cost.

\paragraph{Step 3: Set Up Constraints}
Clearly specify the constraints of the problem. Constraints can be represented using formulas that involve the variables, coefficients, and right-hand side values. For example, a constraint ensuring that total production does not exceed capacity can be written as \texttt{=SUMPRODUCT(Production, Coefficients) <= Capacity}.

\paragraph{Step 4: Open Solver}
Ensure Solver is enabled in Excel. If not, you can activate it by navigating to \texttt{File > Options > Add-Ins > Manage Excel Add-ins} and checking the box for Solver.

Once Solver is enabled, go to \texttt{Data > Solver} to open the Solver Parameters dialog box.
{\centering
\altincludegraphics[width = 0.85\textwidth]{Figures/excel-solver-button}{Screenshot of Excel's Data ribbon with the Solver button highlighted on the right side of the toolbar, showing where to click to launch the Solver Parameters dialog after enabling the add-in.}\par}

\paragraph{Step 5: Define Solver Parameters}
In the Solver Parameters dialog box:
\begin{itemize}
    \item Set the \textbf{Objective} to the cell containing the objective function.
    \item Choose \textbf{Max} or \textbf{Min} to indicate whether the objective is to be maximized or minimized.
    \item Specify the \textbf{Variable Cells} by selecting the range of cells representing decision variables.
    \item Add \textbf{Constraints} by specifying the relationship (e.g., \texttt{<=}, \texttt{>=}, \texttt{=}) and the cells involved.
    \item When the variables are non-negative, you can select \texttt{Make Unconstrained Variables Non-Negative}.
\end{itemize}

{\centering
\altincludegraphics[scale = 0.25]{Figures/excel-solver-solve}{Solver Parameters dialog with fields filled in: Set Objective cell, Max/Min selector, By Changing Variable Cells range, Subject to the Constraints list, and the Make Unconstrained Variables Non-Negative checkbox enabled.}\par}

\paragraph{Step 6: Choose a Solving Method}
Solver provides three solving methods:
\begin{itemize}
    \item \textbf{GRG Nonlinear:} For smooth nonlinear problems.
    \item \textbf{Simplex LP:} For linear programming problems.
    \item \textbf{Evolutionary:} For non-smooth problems or those with integer constraints.
\end{itemize}
Select the method that matches the problem type.
{\centering
\altincludegraphics[scale = 0.3]{Figures/excel-solver-method}{Close-up of the Select a Solving Method dropdown in the Solver Parameters dialog, showing the three options GRG Nonlinear, Simplex LP, and Evolutionary that the user picks based on whether the problem is smooth nonlinear, linear, or non-smooth/integer.}\par}

\paragraph{Step 7: Solve the Problem}
Click \texttt{Solve} to run Solver. 
If Solver finds a solution, it will display the results and allow you to keep or discard them. Analyze the solution to ensure it meets the problem's requirements.

{\centering
\altincludegraphics[scale = 0.32]{Figures/excel-solver-select}{Solver Results dialog that appears after solving, offering Keep Solver Solution or Restore Original Values, with a Reports list (Answer, Sensitivity, Limits) on the right that the user can select to generate diagnostic reports.}\par}

\paragraph{Step 8: Interpret the Results}
Review the values of the decision variables and the objective function. Check that all constraints are satisfied and assess the solution's feasibility and quality.
\begin{itemize}
\item Answer report
{\centering
\altincludegraphics[scale = 0.4]{Figures/excel-solver-answer-report}{Excel Answer Report worksheet generated by Solver, listing the objective cell's original and final values, a table of variable cells with their final values, and a constraints table showing each constraint's cell value, formula, status (Binding or Not Binding), and slack.}\par}
\item Sensitivity report
{\centering
\altincludegraphics[scale = 0.4]{Figures/excel-solver-sensitivity}{Excel Sensitivity Report worksheet with a Variable Cells table (final value, reduced cost, objective coefficient, allowable increase/decrease) and a Constraints table (shadow price, right-hand side, allowable increase/decrease) for each decision variable and constraint.}\par}
\item Solution on spreadsheet
{\centering
\altincludegraphics[scale = 0.4]{Figures/excel-solver-solution}{Excel spreadsheet after Solver has run, with the yellow variable cells now populated with optimal decision-variable values and the green formula cells displaying the resulting objective value and constraint sums, illustrating the final optimized layout.}\par}
\end{itemize}

\begin{resource}
Examples and guides
\begin{itemize}
\item \href{https://www.youtube.com/watch?v=LKV6fT8xApA&ab_channel=JoshuaEmmanuel}{Installing Excel Solver on Windows}
\item \url{https://www.excel-easy.com/data-analysis/solver.html}
\item \href{https://www.youtube.com/watch?v=dRm5MEoA3OI&ab_channel=LeilaGharani}{Excel Solver - Introduction on Youtube}

\item \href{https://ocw.mit.edu/courses/sloan-school-of-management/15-053-optimization-methods-in-management-science-spring-2013/tutorials/MIT15_053S13_tut03.pdf}{Some notes from MIT}
\end{itemize}
Some videos
\begin{itemize}
\item \href{https://www.youtube.com/watch?v=V5DmekIFenA}{Solving a linear program}
\item \href{https://www.youtube.com/watch?v=6xa1x_Iqjzg}{Optimal product mix}
\item \href{https://www.youtube.com/watch?v=vpodCzRtUMU}{Set Cover}

\item \href{https://www.youtube.com/watch?v=tKV24jzZ10s}{Introduction to Designing Optimization Models Using Excel Solver}

\item \href{https://www.youtube.com/watch?v=-E3rSoClgMI}{Traveling Salesman Problem}

\item \href{https://www.youtube.com/watch?v=UQYJvSjXE6I}{Also Traveling Salesman Problem}

\item \href{https://www.youtube.com/watch?v=owHq3Mbniqo}{Multiple Traveling Salesman Problem}

\item \href{https://www.youtube.com/watch?v=JkZkGxVZ8ao}{Shortest Path}

\end{itemize}
Some other links
\begin{itemize}
\item \href{https://github.com/lobodemonte/excel-solver-loan-example}{Loan Example}

\item \href{https://github.com/Bhargavanarasimhan/Excel-Solver-Files}{Several Examples including TSP}
\end{itemize}
\end{resource}

\begin{tryit}
\vizlink{excel-solver}{Excel Solver Walkthrough}: cells, formulas, and the Solver dialog on a small LP.
\end{tryit}


\section{Exercises}

\exgroup{Warm-ups}

\begin{ex}{Solver Rerun: Minimum-Cost Network Flow}{}
\label{ex:excel-network-flow-rerun}
\stars{1} This exercise mirrors Example~\ref{ex:min-cost-flow-warehouses} (Minimum-Cost Network Flow).  Download the Excel workbook linked in that example and change only the cost data: shipping now costs $3$ per unit from Warehouse 1 to Store 1, $7$ from Warehouse 1 to Store 2, $4$ from Warehouse 2 to Store 1, and $2$ from Warehouse 2 to Store 2.  The supplies ($20$ and $30$), demands ($25$ and $25$), and route capacities ($15$, $10$, $20$, $20$) are unchanged.
\begin{enumerate}
\item Re-solve with the Simplex LP method and report the optimal flow on each of the four routes.
\item As a check, the minimum cost is $160$.  Which routes are at capacity?
\end{enumerate}
\exrefs{\S\ref{sec:excel-solver}, Example~\ref{ex:min-cost-flow-warehouses}}
\end{ex}

\begin{ex}{Solver Rerun: Production Planning over 10 Periods}{}
\label{ex:excel-production-rerun}
\stars{1} This exercise mirrors Example~\ref{ex:production-10period} (Production Planning with 10 Periods).  Download the Excel workbook linked in that example.  Keep the production costs $(10, 12, \dots, 28)$, the holding cost of $5$, and the initial inventory $s_0 = 5$, but replace the demands with $(6, 9, 7, 8, 5, 10, 7, 6, 8, 4)$.
\begin{enumerate}
\item Update the demand cells and re-solve with the Simplex LP method.  Report the production plan and the total cost.  As a check, the minimum cost is $1252$.
\item The optimal plan carries no inventory between periods.  Explain why, by comparing the holding cost with the period-to-period increase in production cost.
\end{enumerate}
\exrefs{\S\ref{sec:excel-solver}, Example~\ref{ex:production-10period}}
\end{ex}

\exgroup{Core problems}

\begin{ex}{Diet Problem in Excel}{}
\label{ex:excel-diet}
\stars{2} Consider the following table indicating the nutritional value of different food types.
\begin{center}
\begin{tabular}%{lccccc}
{|p{0.2\linewidth} | p{0.14\linewidth} | p{0.14\linewidth} | p{0.14\linewidth} | p{0.14\linewidth} | p{0.14\linewidth} |}
\hline & Price $(\$)$ per serving & Calories per serving & Fat $(\mathrm{g})$ per serving & Protein $(\mathrm{g})$ per serving & Carbohydrate $(\mathrm{g})$ per serving \\
\hline Raw carrots & $0.14$ & 23 & $0.1$ & $0.6$ & 6 \\
Baked potatoes & $0.12$ & 171 & $0.2$ & $3.7$ & 30 \\
Wheat bread & $0.2$ & 65 & 0 & $2.2$ & 13 \\
Cheddar cheese & $0.75$ & 112 & $9.3$ & 7 & 0 \\
Peanut butter & $0.15$ & 188 & 16 & $7.7$ & 2 \\
\hline
\end{tabular}
\end{center}
You need to decide how many servings of each food to buy each day so that you minimize the total cost of buying your food while satisfying the following daily nutritional requirements:
\begin{itemize}
\item calories must be at least 2000,
\item fat must be at least $50 \mathrm{~g}$,
\item protein must be at least $100 \mathrm{~g}$,
\item carbohydrates must be at least $250 \mathrm{~g}$.
\end{itemize}
Write an LP that will help you decide how many servings of each of the aforementioned foods are needed to meet all the nutritional requirements, while minimizing the total cost of the food (you may buy fractional numbers of servings).

Solve this problem using Excel Solver.  Provide a screenshot of your spreadsheet with the optimal solution.
\exrefs{\S\ref{sec:excel-solver}}
\end{ex}

\begin{ex}{Production Planning in Excel}{}
\label{ex:excel-production-planning}
\stars{2} A small factory produces two products: widgets and gadgets.  Each widget earns \$8 in profit and each gadget earns \$6.  Production is limited by two resources:

\begin{center}
\begin{tabular}{c|c|c|c}
 & \textbf{Widget} & \textbf{Gadget} & \textbf{Available} \\
\hline
Machine time (hrs) & 2 & 1 & 40 \\
Raw material (kg) & 1 & 2 & 30 \\
\end{tabular}
\end{center}

Both products require nonnegative production quantities.

\begin{enumerate}
\item Set up this problem in Excel, clearly color-coding data cells (blue), variable cells (yellow), and formula cells (green).
\item Solve using Excel Solver with the Simplex LP method. Report the optimal production quantities and maximum profit.
\item Generate a Sensitivity Report.  Using the report, answer: if the profit per widget increases to \$9, does the optimal production plan change?
\end{enumerate}
\exrefs{\S\ref{sec:excel-solver}}
\end{ex}

\begin{ex}{Transportation Problem in Excel}{}
\label{ex:excel-transportation}
\stars{2} A company has two warehouses (W1 and W2) that supply three retail stores (S1, S2, S3). The shipping cost per unit, supply at each warehouse, and demand at each store are given below:

\begin{center}
\begin{tabular}{c|c|c|c|c}
 & S1 & S2 & S3 & \textbf{Supply} \\
\hline
W1 & \$4 & \$8 & \$1 & 60 \\
W2 & \$6 & \$3 & \$5 & 50 \\
\hline
\textbf{Demand} & 30 & 40 & 40 & \\
\end{tabular}
\end{center}

\begin{enumerate}
\item Set up a transportation model in Excel with decision variables representing the number of units shipped from each warehouse to each store.
\item Solve using Excel Solver to minimize total shipping cost. Report the optimal shipping plan and minimum cost.
\end{enumerate}
\exrefs{\S\ref{sec:excel-solver}, \S\ref{sec:transportation-problem}}
\end{ex}

\begin{ex}{Product Mix with a Sensitivity Report}{}
\label{ex:excel-furniture-sensitivity}
\stars{2} A furniture shop makes tables, chairs, and desks with profits of \$70, \$50, and \$90 per item.  Weekly resources are 240 hours of carpentry, 100 hours of finishing, and 2000 board-feet of wood:

\begin{center}
\begin{tabular}{c|c|c|c|c}
 & \textbf{Table} & \textbf{Chair} & \textbf{Desk} & \textbf{Available} \\
\hline
Carpentry (hrs) & 4 & 3 & 6 & 240 \\
Finishing (hrs) & 2 & 1 & 3 & 100 \\
Wood (bd-ft) & 30 & 20 & 40 & 2000 \\
\end{tabular}
\end{center}

Production quantities may be fractional.
\begin{enumerate}
\item Set up the LP in Excel with color-coded data, variable, and formula cells, and solve with the Simplex LP method.  Report the production plan and the maximum profit.  As a check, the maximum profit is \$4100.
\item Generate a Sensitivity Report.  Report the shadow price of each resource and state which constraints are binding.
\item Using only the report, decide: is it worth paying \$10 per hour for additional finishing time?
\item The report lists an allowable increase for the desk profit coefficient.  If the profit per desk rises to \$100, does the production plan change?  What if it rises to \$110?
\end{enumerate}
\exrefs{\S\ref{sec:excel-solver}}
\end{ex}

\begin{ex}{Unbalanced Transportation with a Sensitivity Report}{}
\label{ex:excel-unbalanced-transportation}
\stars{2} Three plants supply three cities.  Plant capacities are 70, 50, and 40 units; city demands are 50, 60, and 40 units, so total supply exceeds total demand.  Shipping costs per unit are:

\begin{center}
\begin{tabular}{c|c|c|c|c}
 & C1 & C2 & C3 & \textbf{Supply} \\
\hline
P1 & \$4 & \$5 & \$7 & 70 \\
P2 & \$6 & \$3 & \$2 & 50 \\
P3 & \$5 & \$8 & \$4 & 40 \\
\hline
\textbf{Demand} & 50 & 60 & 40 & \\
\end{tabular}
\end{center}

\begin{enumerate}
\item Set up the model in Excel with supply constraints as $\leq$ and demand constraints as $=$, and solve with the Simplex LP method.  Report the shipping plan and the minimum cost.  As a check, the minimum cost is \$560.
\item Which plant does not ship its full capacity, and how many units stay there?
\item Generate a Sensitivity Report.  Report the shadow price of each supply constraint and interpret the sign of the nonzero one: what would one extra unit of capacity at that plant be worth?
\end{enumerate}
\exrefs{\S\ref{sec:excel-solver}, \S\ref{sec:transportation-problem}}
\end{ex}

\exgroup{Concepts and connections}

\begin{ex}{Simplex LP versus GRG Nonlinear}{}
\label{ex:excel-solver-methods}
\stars{2} Solver's dialog offers three solving methods: Simplex LP, GRG Nonlinear, and Evolutionary.
\begin{enumerate}
\item Describe what class of problems each of Simplex LP and GRG Nonlinear is designed for.
\item If your model is a linear program, why should you choose Simplex LP rather than GRG Nonlinear, even though both may return the same answer?  Consider the guarantee on the solution (global versus possibly local optimum) and the contents of the Sensitivity Report.
\item Suppose you select Simplex LP but one of your constraint cells contains the formula \texttt{=B2*B3}, where both \texttt{B2} and \texttt{B3} are variable cells.  What will Solver do, and why?
\item Give an example of an objective function for which GRG Nonlinear is appropriate and Simplex LP is not.
\end{enumerate}
\exrefs{\S\ref{sec:excel-solver}}
\end{ex}

\exgroup{Challenge problems}

\begin{ex}{When a Variable Must Go Negative}{}
\label{ex:excel-short-selling}
\stars{3} An investor has \$100 (in thousands) to split between fund A, which returns 5\%, and fund B, which returns 12\%.  The broker allows \emph{short selling} fund A: the investor may hold a negative amount of A and use the proceeds to buy more of B.  Margin rules cap the holding in B at \$140.  The model is
\begin{align*}
\max \quad & 0.05\,x_A + 0.12\,x_B\\
\text{s.t.} \quad & x_A + x_B = 100\\
& x_B \leq 140, \qquad x_B \geq 0, \qquad x_A \text{ free}.
\end{align*}
\begin{enumerate}
\item Set up the model in Excel and solve with the Simplex LP method, leaving \texttt{Make Unconstrained Variables Non-Negative} checked.  Report the solution.
\item Now uncheck \texttt{Make Unconstrained Variables Non-Negative}, add the constraint $x_B \geq 0$ explicitly, and re-solve.  Report the new solution.
\item Explain exactly what the checkbox does and why it must be unchecked for this model.  What risk do you take on for other variables when you uncheck it?
\end{enumerate}
\exrefs{\S\ref{sec:excel-solver}}
\end{ex}

\subsubsection*{Selected Solutions}

\begin{solution}(Exercise~\ref{ex:excel-production-planning})
Let $x_1$ and $x_2$ be the number of widgets and gadgets produced.  The LP behind the spreadsheet is
\begin{align*}
\max \quad & z = 8x_1 + 6x_2\\
\text{s.t.} \quad & 2x_1 + x_2 \leq 40 && \text{(machine time)}\\
& x_1 + 2x_2 \leq 30 && \text{(raw material)}\\
& x_1, x_2 \geq 0.
\end{align*}
In Excel, put the two variables in yellow cells, compute the objective and the two constraint left-hand sides with \texttt{SUMPRODUCT} in green cells, and add the constraints in Solver with the Simplex LP method.  Solver returns $x_1 = 50/3 \approx 16.67$ widgets and $x_2 = 20/3 \approx 6.67$ gadgets with maximum profit $520/3 \approx \$173.33$; both resource constraints are binding.
For part 3, the Sensitivity Report shows the allowable increase for the widget profit coefficient.  Raising it from \$8 to \$9 stays within the allowable range, so the optimal production plan does not change --- only the profit does, rising to \$190.  (Geometrically, the objective slope stays between the slopes of the two binding constraints, so the same corner point remains optimal.)
\end{solution}

\begin{solution}(Exercise~\ref{ex:excel-transportation})
Let $x_{ij} \geq 0$ be the number of units shipped from warehouse $i$ to store $j$.  A natural Excel layout is a $2 \times 3$ grid of yellow variable cells mirroring the cost table, with row sums compared to supply and column sums compared to demand:
\begin{align*}
\min \quad & z = 4x_{11} + 8x_{12} + x_{13} + 6x_{21} + 3x_{22} + 5x_{23}\\
\text{s.t.} \quad & x_{11} + x_{12} + x_{13} \leq 60 && \text{(supply W1)}\\
& x_{21} + x_{22} + x_{23} \leq 50 && \text{(supply W2)}\\
& x_{11} + x_{21} = 30, \quad x_{12} + x_{22} = 40, \quad x_{13} + x_{23} = 40 && \text{(demands)}\\
& x_{ij} \geq 0.
\end{align*}
Total supply equals total demand (110 units), so both supplies are used fully.  Solver finds the optimal plan $x_{11} = 20$, $x_{13} = 40$, $x_{21} = 10$, $x_{22} = 40$ (all other routes unused), with minimum cost \$300: each warehouse ships everything it can along its cheapest routes (W1 to S3, W2 to S2), and store S1's demand is split to make the supplies balance.
\end{solution}

\begin{solution}(Exercise~\ref{ex:excel-short-selling})
With the checkbox checked, Solver silently adds $x_A \geq 0$ and $x_B \geq 0$, so the best it can do is $x_A = 0$, $x_B = 100$, for a return of $0.05 \cdot 0 + 0.12 \cdot 100 = 12$ (that is, \$12{,}000).  With the checkbox unchecked and $x_B \geq 0$ added as an explicit constraint, Solver finds $x_A = -40$, $x_B = 140$: the investor shorts \$40 of fund A to push fund B to its margin cap, earning $0.05(-40) + 0.12(140) = 14.8$, or \$14{,}800.
The checkbox applies a lower bound of zero to every decision variable that has no explicit lower bound in the constraint list.  Short selling is exactly a negative holding, so that hidden bound cuts off the optimal solution.  The risk when unchecking is that all other variables also lose the automatic bound: any variable that should be non-negative (here $x_B$) must now get its own $\geq 0$ constraint, or Solver may return meaningless negative values for it.
\end{solution}

